% Part: first-order-logic
% Chapter: introduction
% Section: substitution

\documentclass[../../../include/open-logic-section]{subfiles}

\begin{document}

\olfileid[fr]{fol}{int}{sub}

\olsection{Substitution}

Examinons un exemple qui montrera comment les différentes notions
s'articulent et comment l'élaboration de la syntaxe et de la sémantique
prépare les études plus avancées qui suivront. Nos systèmes de
!!{derivation} doivent nous permettre de !!{derive}
$\Atom{\Obj P}{\Obj a}$ à partir de
$\lforall[\Obj v_0][\Atom{\Obj P}{\Obj v_0}]$. Nous pourrions même
vouloir présenter cela comme une règle d'inférence. Pour le faire, il
faut toutefois l'énoncer dans toute sa généralité : non pas seulement
pour $\Obj P$, $\Obj a$ et $\Obj v_0$, mais pour une
!!{formula}~$!A$, un terme~$t$ et une !!{variable}~$x$ quelconques.
Rappelons que les !!{constant}s sont des termes; nous considérerons
aussi des termes plus complexes construits à partir de !!{constant}s
et de !!{function}s. Nous voulons donc pouvoir dire quelque chose
comme : « Chaque fois que l'on a dérivé
$\lforall[x][!A(x)]$, on est fondé à inférer~$!A(t)$, résultat obtenu
en supprimant $\lforall[x]$ et en remplaçant~$x$ par~$t$. » Mais que
signifie exactement « remplacer $x$ par~$t$ » ? Quelle relation y
a-t-il entre $!A(x)$ et~$!A(t)$ ? Cette opération est-elle toujours
licite ?

Pour rendre cette idée précise, nous définissons l'opération de
\emph{substitution}. La substitution est en réalité délicate : nous ne
pouvons pas simplement remplacer tous les~$x$ de~$!A$ par~$t$, et
on ne peut pas toujours substituer n'importe quel terme~$t$ à
n'importe quelle variable~$x$. Nous traiterons encore ces difficultés
au moyen de définitions inductives. Une fois ces définitions données,
la règle « inférer $!A(t)$ de $\lforall[x][!A(x)]$ » sera définie avec
précision. Nous pourrons en outre montrer que cette règle d'inférence
est correcte, au sens où $\lforall[x][!A(x)]$ entraîne~$!A(t)$. Pour le
démontrer, il faudra encore établir un résultat dont l'utilité peut
d'abord sembler obscure. Le fait que $\lforall[x][!A(x)]$
entraîne~$!A(t)$ repose sur la propriété suivante : savoir si
$\Sat{M}{!A(t)}$ vaut ne dépend que de la valeur du terme~$t$. Si $m$
est l'!!{element} de~$\Domain{M}$
désigné par~$t$, alors
$\Sat{M}{!A(t)}[s]$ si et seulement si
$\Sat{M}{!A(x)}[\Subst{s}{m}{x}]$. Ce résultat reste vrai lorsque $t$
contient des !!{variable}s, mais il faudra l'énoncer avec soin.

\begin{editorial}
Dans la première formule de la source, les arguments de
$\Atom{\Obj P}{\Obj v_0}$ et de
$\lforall[\Obj v_0][\Atom{\Obj P}{\Obj v_0}]$ sont délimités de façon
incompatible avec les macros employées. Les limites correspondant à la
formule universelle décrite par la phrase sont rétablies.
\end{editorial}

\end{document}
